% =====================================================================
\chapter{Funciones, imagen y preimagen}
\label{tema:2}
% =====================================================================

Este tema define con precisión qué es una función y estudia cómo actúa sobre subconjuntos mediante la imagen y la preimagen. Sigue las secciones~12.1 y~12.6 de Hammack~\cite{hammack}, las secciones~5.1 y~5.5 de Velleman~\cite{velleman} y la sección~4.3 de~\cite{mcs}. Las propiedades de la imagen y la preimagen son las de la pregunta~9 de la hoja~1 de Cambridge~\cite{camns}.

\section{Definición de función}

\begin{definicion}[Función]\label{def:t2-funcion}
Sean $A$ y $B$ conjuntos. Una \emph{función} $f$ de $A$ en $B$, escrita $f\colon A\to B$, es un subconjunto $f\subseteq A\times B$ tal que para cada $a\in A$ existe \emph{exactamente un} $b\in B$ con $(a,b)\in f$; ese $b$ se denota $f(a)$. El conjunto $A$ es el \emph{dominio} y $B$ el \emph{codominio}. El \emph{recorrido} (o \emph{imagen}) de $f$ es $f(A)=\{f(a): a\in A\}\subseteq B$.
\end{definicion}

La definición contiene dos exigencias: \emph{existencia} (todo elemento del dominio tiene imagen) y \emph{unicidad} (no tiene dos). Dos funciones $f,g\colon A\to B$ son iguales si $f(a)=g(a)$ para todo $a\in A$; la igualdad exige además que el dominio y el codominio coincidan.

\addtocounter{definicion}{1}%

\addtocounter{definicion}{1}%

\section{El conjunto de funciones}

\begin{definicion}
$B^{A}$ denota el conjunto de todas las funciones $A\to B$.
\end{definicion}

\begin{proposicion}\label{prop:t2-contar}
Si $A$ y $B$ son finitos, $\abs{B^{A}}=\abs{B}^{\abs{A}}$.
\end{proposicion}

\begin{definicion}[Función característica]
Sea $X\subseteq U$. La \emph{función característica} (o \emph{indicadora}) de $X$ es $\chi_X\colon U\to\{0,1\}$, con $\chi_X(u)=1$ si $u\in X$ y $\chi_X(u)=0$ si $u\notin X$.
\end{definicion}

Cada subconjunto de $U$ determina su función característica y cada función $U\to\{0,1\}$ es la característica de un único subconjunto, el de los elementos que van al~1. Con $\abs{U}=n$, la \cref{prop:t2-contar} da $\abs{\{0,1\}^{U}}=2^{n}$, de nuevo el \cref{teo:t1-partes}. En el \cref{tema:3} se formaliza esta correspondencia como una biyección.

\section{Imagen y preimagen de subconjuntos}

\begin{definicion}[Imagen y preimagen]\label{def:t2-imagen}
Sea $f\colon A\to B$. Para $X\subseteq A$ y $Y\subseteq B$ se definen la \emph{imagen} de $X$ y la \emph{preimagen} (o \emph{imagen inversa}) de $Y$:
\[
f(X)=\{f(x): x\in X\}\subseteq B,\qquad f^{-1}(Y)=\{a\in A: f(a)\in Y\}\subseteq A.
\]
\end{definicion}

La notación $f^{-1}(Y)$ tiene sentido para \emph{cualquier} función, sea o no invertible: es un conjunto, no la aplicación de una función inversa. Si $Y=\{b\}$, el conjunto $f^{-1}(\{b\})$ se llama \emph{fibra} de $b$.

\addtocounter{definicion}{1}%

\begin{teorema}[Propiedades de la preimagen]\label{teo:t2-preimagen}
Sea $f\colon A\to B$ y sean $C,D\subseteq B$. Entonces
\begin{gather*}
f^{-1}(C\cup D)=f^{-1}(C)\cup f^{-1}(D),\qquad f^{-1}(C\cap D)=f^{-1}(C)\cap f^{-1}(D),\\
f^{-1}(B\setminus C)=A\setminus f^{-1}(C).
\end{gather*}
\end{teorema}

\begin{teorema}[Propiedades de la imagen]\label{teo:t2-imagen}
Sea $f\colon A\to B$, $X_1,X_2,X\subseteq A$ e $Y\subseteq B$. Entonces
\begin{enumerate}
  \item $f(X_1\cup X_2)=f(X_1)\cup f(X_2)$;
  \item $f(X_1\cap X_2)\subseteq f(X_1)\cap f(X_2)$, y la inclusión puede ser estricta;
  \item $X\subseteq f^{-1}(f(X))$ y $f(f^{-1}(Y))\subseteq Y$.
\end{enumerate}
\end{teorema}

La asimetría entre ambos teoremas es instructiva: la preimagen respeta todas las operaciones de conjuntos, mientras que la imagen solo respeta la unión. La razón es que la imagen puede «pegar» elementos distintos; en el \cref{tema:3} se verá que las igualdades $f(X_1\cap X_2)=f(X_1)\cap f(X_2)$ y $f^{-1}(f(X))=X$ son ciertas exactamente cuando $f$ es inyectiva.

\section{Otras construcciones}

Una \emph{sucesión} de elementos de $A$ es una función $\N\to A$, $n\mapsto a_n$. Una \emph{cadena} de longitud $n$ sobre $\Sigma$ puede verse como una función $\{1,\dots,n\}\to\Sigma$ que asigna a cada posición su símbolo. Si $f\colon A\to B$ y $X\subseteq A$, la \emph{restricción} $f|_X\colon X\to B$ es la función $f|_X(x)=f(x)$. Las funciones \emph{definidas a trozos} asignan reglas distintas a partes disjuntas del dominio que lo recubren.

\addtocounter{definicion}{1}%

\addtocounter{definicion}{1}%

